% Section 2. Soft-sensor reframe 2026-11-04 (branch docs/cep-soft-sensor-frame), from hardware.tex (sec:plant) and
% csd.tex (sec:expected, tab:microscope) of docs/cep-grounded-draft. Numbers: scripts/microscope_summary.py,
% scripts/check_anchors.py; dataset facts: datasets/README.md. Claim changes: papers/cep/CLAIM_CHANGES.md.
\section{Process and data}\label{sec:process}

This section gives the process, the recordings, the process signals and the microscope reference that the
measurement model (RQ1) and the estimator (RQ3) use.

\subsection{The seed pan and its strikes}\label{sec:plant}

The camera is installed on the \textbf{seed pan} of the Slovensk\'e cukrovary sugar plant in Sere\v{d}, Slovakia: the
vacuum pan designated ZRA in the plant's DCS (Slovak \emph{o\v{c}kovac\'i varostroj}; plant Q\&A, 2026-08-12, and the
DeltaV ``Udaje ZRA'' screen, \citealp{SlovenskeCukrovary2026}; English--Slovak glossary in
\ref{app:glossary}). A strike starts from syrup, which is concentrated to supersaturation and seeded with seed
slurry (\emph{mikroo\v{c}ko}), a suspension of milled sucrose in a non-solvent liquid. The crystals then grow during
boil-up while syrup is fed, and the massecuite is tightened (\emph{su\v{s}enie}) and discharged. The product is seed
massecuite (\emph{o\v{c}kovacia cukrovina}), which is dosed into the A pans (operator A.~Fedor, pers.\ comm.,
2026-11-03). The pan is mechanically stirred in normal operation (stirrer current tag EI1004), which is the basis
for treating its contents as a well-mixed massecuite (\secref{sec:layer}).

The camera views the pan contents through a circular sight glass in a bolted flange, about 150~mm in diameter and
20~mm thick (plant statements). A lamp inside the pan, an optical fibre ``very close to the glass'' (plant statement;
an estimated 1--2~mm from the inner face, not measured), backlights the massecuite, so crystals appear as dark
outlines on a bright field.

\subsection{Recordings, process signals and samples}\label{sec:data}

\tabref{tab:data} lists the data. The \emph{sampled run of 2026-10-30} covers two strikes, b1 and b2, seeded at
11:29:33 and 15:22:00 (run sidecar). Its camera recording is a lossless $720 \times 576$, 25~fps file; its capture
path (recorder export or capture card) is not recorded \pending{capture chain of the 2026-10-30 recording}.
Samples were drawn from the pan during both strikes. Each was photographed under a microscope (\secref{sec:microscope})
and its refractometric dry substance and polarisation were measured, from which Brix and purity follow (run
sidecar). No DCS export exists for this day \pending{DCS export for 2026-10-30, requested from the plant}.

The \emph{PAL recording of 2026-02-04} is the recorder's own $720 \times 576$ file. It covers strike b2 from before
seeding (12:30, operator note; steam valve cut at 12:31:30 in the DCS) to discharge (16:00:10, the level passing the
glass). A DCS export of the seed pan for the same afternoon, at 30~s, holds the heating-steam pressure (PIC1001), the
pan vacuum (PIC1002), the signal QI1003, the syrup-feed controller (QIC1000), the massecuite temperature (TI1006),
the stirrer current (EI1004) and the receiver level (LI1009), but not the pan level LI1000. QI1003 is labelled
viscosity in the DCS and reads on a Brix scale; it measures the whole massecuite, not the mother liquor, so it is not
a true Brix. The recorder clock agrees with the DCS clock to about $\pm 30$~s (\code{scripts/check_anchors.py}). No
size reference exists for this day. Three 1080p screen recordings of disputed date have no size reference either;
they are used only in \ref{app:detectors}.

Every frame is cropped to the live camera picture before analysis: the recorder's date band and channel label are
sharp digital glyphs that would otherwise dominate any sharpness statistic (\ref{app:optics}).

\begin{table*}
\centering\tabfont
\caption{Data used in this paper (\texttt{datasets/README.md}). DCS: distributed control system export.}\label{tab:data}
\begin{tabular}{@{}>{\raggedright\arraybackslash}p{0.17\textwidth}>{\raggedright\arraybackslash}p{0.2\textwidth}%
>{\raggedright\arraybackslash}p{0.22\textwidth}>{\raggedright\arraybackslash}p{0.12\textwidth}%
>{\raggedright\arraybackslash}p{0.21\textwidth}@{}}
\toprule
Run & Camera & Size reference & DCS & Use \\
\midrule
2026-10-30, sampled run, strikes b1 and b2 & lossless $720 \times 576$, 25~fps & microscope photographs of samples
(\tabref{tab:microscope}); lab Brix and purity & none (requested) & measurement model, anchor, estimator validation \\
2026-02-04, strike b2 & recorder PAL file, $720 \times 576$ & none & 30~s, no LI1000 & failed classical-detector baseline
(\secref{sec:csd-time}); baseline inputs \\
1080p clips, date disputed & screen recordings of the recorder output & none & none & \ref{app:detectors} only \\
Seed slurry, labelled 2025-01-16 & -- & laser diffraction \citep{Mastersizer2025} & -- & size at seeding \\
\bottomrule
\end{tabular}
\end{table*}

\subsection{Microscope reference}\label{sec:microscope}

The seed slurry dosed into the seed pan has $\mathrm{Dv}(10)/\mathrm{Dv}(50)/\mathrm{Dv}(90) = 1.57/6.77/18.6$~\um{}
(laser diffraction in 2-propanol, mean of three runs, measured 2025-03-04 by the AGRANA central laboratory;
\citealp{Mastersizer2025}). In the sampled run of 2026-10-30, crystals in the samples were measured under a
microscope: one chord per operator-picked crystal, burned into 54 photographs and transcribed into
\code{datasets/microscope_annotations.csv} (1052 readable values, one AI reading; a human second reading is
\pending{second reading of the quoted samples}; \tabref{tab:microscope}). The microscope camera clock is known to
$\pm 47$~s. The seed-slurry photograph shows 10--19~\um. The photograph taken just before seeding in b1 already
shows 18--42~\um{} crystals (median 33~\um), larger than the 29~\um{} median 2~min after seeding; whether they are
residual massecuite or spontaneous nucleation is open. From about 30--40~min the median exceeds 0.1~mm, and 122~min
after seeding it is 0.48~mm (b1, range 0.22--1.14~mm). At equal times the strikes differ: at 96~min the b1 median is
355~\um{} and the b2 median 225~\um. These are chords of operator-picked crystals, not distributions; the largest
single values may be agglomerates, and the 2--3~\um{} chords at the b2 seeding are a few microscope pixels long,
below the optical resolution. We define the crystal size $L$ of this paper as the median microscope chord, the
laboratory definition against which the camera is calibrated. \secref{sec:measurement} relates the camera signal to
it.

\begin{table}
\centering\tabfont
\caption{Microscope chords of operator-picked crystals against time since seeding $t$, seed pan, sampled run of
2026-10-30 (photographs, single AI transcription; \texttt{scripts/\allowbreak microscope\_summary.py}). Selected
samples; all 20 are in \texttt{datasets/\allowbreak README.md}. $^\dagger$Taken 27~s before seeding.
$^\ddagger$Includes 2--3~\um{} chords below the optical resolution.}\label{tab:microscope}
\begin{tabular}{@{}llrrr@{}}
\toprule
Strike & Sample & $t$ (min) & $n$ & min / median / max (\textmu m) \\
\midrule
-- & seed slurry & -- & 4 & 10 / 13 / 19 \\
b1 & 11:30$^\dagger$ & 0 & 6 & 18 / 33 / 42 \\
b1 & 11:32 & 2 & 13 & 22 / 29 / 35 \\
b1 & 12:07 & 37 & 50 & 38 / 114 / 271 \\
b1 & 12:30 & 60 & 37 & 62 / 199 / 317 \\
b1 & 13:06 & 96 & 35 & 154 / 355 / 596 \\
b1 & 13:32 & 122 & 26 & 220 / 483 / 1137 \\
b1 & 14:20 & 170 & 51 & 189 / 452 / 995 \\
b2 & 15:22$^\ddagger$ & 0 & 20 & 2 / 10 / 23 \\
b2 & 15:52 & 30 & 50 & 22 / 103 / 218 \\
b2 & 16:31 & 69 & 121 & 40 / 209 / 628 \\
b2 & 16:58 & 96 & 109 & 87 / 225 / 808 \\
b2 & 17:48 & 146 & 103 & 117 / 299 / 695 \\
b2 & 18:04 & 162 & 65 & 101 / 399 / 961 \\
\bottomrule
\end{tabular}
\end{table}
